% Generated by tools/edn_latex.py from reports/edn.md.
% Do not edit: change reports/edn.md and run `make edn`.
\ednentry{
  id = {EDN-19},
  title = {Which DagsHub repository the team uses as DVC remote and MLflow server},
  date = {2026-09-29},
  milestone = {M2: Reproducibility},
  topic = {Data Versioning, Experiment Tracking},
  participants = {@lukas2510 (Scrum Master), on behalf of the team. Pending confirmation by the full team at the next sprint review.},
  decision = {A DagsHub \textbf{organisation} owned by the team hosts the connected repository, and every team member joins it with push rights (option A).
The repository is \href{https://dagshub.com/recommenditos/MLOPS_Recommenditos}{\texttt{https:/\allowbreak{}/\allowbreak{}dagshub.\allowbreak{}com/\allowbreak{}recommenditos/\allowbreak{}MLOPS\_\allowbreak{}Recommenditos}}, owned by the \texttt{recommenditos} DagsHub organisation and connected through DagsHub's GitHub integration.
A first decision for \texttt{pauadal03/\allowbreak{}MLOPS\_\allowbreak{}Recommenditos} was taken and withdrawn on the same day, once it turned out that pauadal03 is not a member of this team: they hold read-only access to our GitHub repository, have never contributed to it, and are not on the roster.
Their DagsHub repository is a mirror of our public repository made from outside the team, so it is not a candidate.
Later the same day \texttt{mark.\allowbreak{}welf.\allowbreak{}atzberger/\allowbreak{}MLOPS\_\allowbreak{}Recommenditos} was deleted, so that candidate no longer exists either and the raw object it held is gone with it.
The practical choice is therefore between a DagsHub organisation owned by the team and a fresh repository under one team member's account.
Whether the repository stays public was decided separately, see EDN-20.},
  rationale = {\emph{Alternatives considered:}
\begin{itemize}[leftmargin=*, topsep=2pt]
\item \textbf{Option A (chosen): a DagsHub organisation owns the connected repository, every team member is a member with push rights.}
\newline Pros: nothing the project depends on hangs off one person's private account; push rights follow membership instead of manual collaborator entries, and the team is five people, so inviting everyone is trivial; visibility is set in one place; data and experiment history survive if a member leaves or cleans up their account.
\newline Cons: one-time setup (create the organisation, connect the GitHub repository, invite four people); under a \texttt{dvc add} mechanism the 548.6 MB raw object has to be pushed once more, though under \texttt{import-\allowbreak{}url} there is nothing to move.
\item \textbf{Option B: use \texttt{mark.\allowbreak{}welf.\allowbreak{}atzberger/\allowbreak{}MLOPS\_\allowbreak{}Recommenditos}, reconnected properly.} No longer available: the repository was deleted on 2026-09-29.
Pros (while it existed): it belonged to a teammate and already held the raw object; no organisation to create.
\newline Cons: its mirror was demonstrably not in sync (see the evidence below), and the likely fix was to delete and re-import it through the GitHub integration, which would have lost the pushed object anyway; the project would still have hung off one teammate's private account, with a manual collaborator entry needed per member.
\item \textbf{Option C: a fresh repository under one team member's account, connected through the GitHub integration.}
\newline Pros: one form less than an organisation; whoever creates it is admin and can add the other four without waiting on anyone.
\newline Cons: identical to option B's structural drawbacks -- the project hangs off one private account and rights are granted per person by hand. The saving over option A is a single setup form.
\item \textbf{Option D: use \texttt{pauadal03/\allowbreak{}MLOPS\_\allowbreak{}Recommenditos}.} Chosen on 2026-09-29 and withdrawn the same day.
\newline Pros: it is the only mirror verifiably connected through the GitHub integration, and adopting it would have cost nothing.
\newline Cons: ruled out on ownership, not on measurements -- pauadal03 is not on the team, has \texttt{read} permission on our GitHub repository and has never contributed to it, so the project's data and experiment history would sit on the account of someone who cannot even push to the repository they mirrored.
For the record, no project data was ever on that mirror: the raw object returns 404 there, so it holds our public git history and nothing else.
\item \textbf{Option E: keep several mirrors and let everyone push to their own.}
\newline Pros: no coordination needed.
\newline Cons: not viable -- the remote URL is committed in \texttt{.\allowbreak{}dvc/\allowbreak{}config}, so the team would overwrite each other's setting, and data and experiment history would be split across accounts.
\end{itemize}
\par
\emph{Why:} By the time the decision was taken the alternatives had collapsed: option D was disqualified on ownership and option B ceased to exist when that repository was deleted, so every remaining candidate required a fresh connect through the GitHub integration anyway.
That removed the only real argument against option A, which had been that a correctly connected repository already existed elsewhere at zero cost.
What was left is a one-form difference -- an organisation versus a repository under one member's account, for the same connect and the same invitations -- against the fact that fifteen cohort members outside the team hold read access to the GitHub repository and can mirror it at any time, as one of them already had.
The team's standing rule was to take the better-engineered option unless it is overkill for a one-semester course project; with four teammates to invite rather than the whole twenty-one-person collaborator list an earlier miscount had assumed, option A is not overkill.
Evidence gathered on 2026-09-29 against the DagsHub and GitHub APIs.
Ownership: write access to the GitHub repository is held by @lukas2510, @kadameit, @ulasawczuk, @W11W11W11 and @michudud04, plus @martinezmatias and @santidrj, the two lecturers who own the course organisation; @pauadal03 has \texttt{read}, zero pull requests and zero commits, and does not appear on the roster in \texttt{docs/\allowbreak{}docs/\allowbreak{}scrum/\allowbreak{}index.\allowbreak{}md}.
Sync: \texttt{pauadal03/\allowbreak{}MLOPS\_\allowbreak{}Recommenditos} is at \texttt{2e902e2143}, identical to GitHub \texttt{main}, carries all 7 branches and mirrors all 7 open issues, while \texttt{mark.\allowbreak{}welf.\allowbreak{}atzberger/\allowbreak{}MLOPS\_\allowbreak{}Recommenditos} is at \texttt{8b1d64fc}, is missing the \texttt{docs/\allowbreak{}data-\allowbreak{}facts-\allowbreak{}corrections} and \texttt{model-\allowbreak{}card} branches, still carries \texttt{feature/\allowbreak{}ruff-\allowbreak{}pl-\allowbreak{}and-\allowbreak{}coverage} which GitHub deleted after the merge, and mirrors 0 issues.
DagsHub only mirrors issues and pull requests for repositories connected through the GitHub integration, so the issue count is the clearest signal that the two were connected in different ways and that mark's is a plain git mirror.
Visibility: both report \texttt{private: false}; EDN-07 as first written on the issue \#3 branch assumed the remote was private, a premise found false and recorded in EDN-20.
The reason that original check misled us: anonymous access is not evidence either way, because DagsHub refuses every anonymous request, including for a known-public control repository (\texttt{DAGsHub-\allowbreak{}Official/\allowbreak{}dagshub-\allowbreak{}docs}, also \texttt{private: false}) and for a repository that does not exist.
Access: \texttt{lukas2510} was granted \texttt{push} on mark's repository on 2026-09-29, shortly before it was deleted, and has \texttt{pull} only on pauadal03's; the rest of the team had push on neither.
State at the end of 2026-09-29: \texttt{mark.\allowbreak{}welf.\allowbreak{}atzberger/\allowbreak{}MLOPS\_\allowbreak{}Recommenditos} returns 404 and so does the raw object it held; \texttt{pauadal03/\allowbreak{}MLOPS\_\allowbreak{}Recommenditos} still exists and still holds no data; a repository search returns that mirror as the only one left.
Structural point behind all of this: fifteen cohort members outside the team hold \texttt{read} on our GitHub repository, so anyone of them can create a mirror at any time without doing anything wrong. The rule the team needs is therefore not ``check who owns a mirror'' but ``the project's infrastructure lives somewhere the team owns'', which is an argument for option A independent of the measurements.
Follow-up: \href{https://github.com/mlops-2627q1-mds-upc/MLOPS_Recommenditos/pull/26}{PR \#26} now commits the organisation's URL in \texttt{.\allowbreak{}dvc/\allowbreak{}config}, and the owner still has to grant push rights to everyone.
Nothing has to be migrated: the 548.6 MB raw object that had been pushed to mark's repository is gone with that repository, so under the \texttt{dvc add} mechanism settled in EDN-25 (\href{https://github.com/mlops-2627q1-mds-upc/MLOPS_Recommenditos/pull/27}{PR \#27}) it is re-downloaded from Zenodo and pushed once to the organisation's remote.
\par
Amended on 2026-10-01, confirmed by the repository owner: all five team members now hold \textbf{write} access to the DagsHub repository, and it is held as a per-repository collaborator entry for each of them rather than as membership of the \texttt{recommenditos} organisation.
The follow-up above is therefore done, by a route this entry did not choose.
The decision is not revisited, because for a single repository the two forms are functionally equivalent: everyone can push, and the repository still belongs to the organisation rather than to a private account, which is what option A was for.
What the deviation does cost is the one property the entry gave as the reason to prefer membership -- rights would follow the person instead of being granted per repository -- so a second repository under the organisation needs its own five grants, and a member who leaves has to be removed once per repository rather than once.
Recorded rather than left to be discovered, so that whoever adds the next repository does not assume the invitations are already in place.},
  ai = {Information seeking, Alternative generation, Alternative assessment, Recommendation},
  response = {Accepted with modifications},
  assessment = {AI found the two mirrors, measured them against GitHub and recommended \texttt{pauadal03}'s on those measurements. The measurements were correct and the recommendation was wrong, because AI never checked whether that repository belonged to anyone on the team -- it ranked the candidates it had found without asking where they came from. Lukas caught it by asking who pauadal03 is, which invalidated the recommendation and the decision taken on it. Recorded here rather than silently corrected, because the failure is the instructive part: a candidate that scores best on every technical measurement can still be disqualified by a question nobody asked. AI's useful contributions were the measurements, the finding that mark's mirror is not properly connected, and the control-repository method behind EDN-20; its first recommendation was not one of them. Its second recommendation, option A, was the one the team accepted, and only after Lukas had rejected the first and asked who the account behind it belonged to.
This is the second decision in the same week reached on a premise nobody had checked, after ``the DagsHub repository is private'' in the first version of EDN-07 on the issue \#3 branch, and in both cases the check that settled it took one API call and was available before the decision rather than after it. The pattern, not either individual fix, is the thing worth carrying into the working agreements: state the premise a decision rests on, and verify it, before recording the decision.},
  aievidence = {Claude Code session on 2026-09-29, prompts translated from German: the comparison and the recommendation for pauadal03 followed ``isn't this our DagsHub repo, the one that matches our GitHub repo?'' and ``why did that repo score worst?''; the retraction followed ``I think pauadal03 is one of our lecturers -- is it bad that everything now has to go through them?''.},
  evidence = {\href{https://github.com/mlops-2627q1-mds-upc/MLOPS_Recommenditos/pull/26}{PR \#26}, \href{https://github.com/mlops-2627q1-mds-upc/MLOPS_Recommenditos/issues/10}{issue \#10}, EDN-20, EDN-25, roster in \texttt{docs/\allowbreak{}docs/\allowbreak{}scrum/\allowbreak{}index.\allowbreak{}md}.},
}
